\section{Conclusion}

We set out to test whether community breadth at benchmark introduction predicts displacement timing. After validating all 5 FAIL FAST gates (G0=0.862, G1=0.605, G2=0.975, G3=0.246, G4=2.14), Cox regression on 258 complete-case rows (EPV $\approx$ 86) yields:

\begin{center}
\textbf{HR = 1.006 (95\% CI = [0.846, 1.196]), LRT $p$ = 0.9495}
\end{center}

This is a near-perfect null ($\Delta\log L = 0.0020$). Neither the lock-in nor the saturation mechanism is operative at a detectable level. Community breadth diversity, as measured by paper submission count, is ruled out as a predictor class for benchmark displacement timing in the Papers With Code ecosystem (2015--2023). This result is specific to the submission-count operationalization; score velocity and institutional diversity remain untested at full power.

The 5-gate FAIL FAST protocol is a reusable methodological contribution for benchmark lifecycle survival analysis --- a pre-registration-equivalent rigor framework that separates data quality failure from hypothesis failure.

Future directions: (1) score trajectory ($\Delta$score\_lag1\_z) as primary predictor on the full 258-row panel; (2) institution-deduplicated breadth via Semantic Scholar author API; (3) cross-platform replication on Semantic Scholar or OpenReview data.
